% =====================================================================
%  PRÓLOGO
% =====================================================================
\capituloSinNumero{Prólogo}

El prólogo presenta, en un tono más personal que el resto del trabajo, las
circunstancias que dieron origen a la monografía: el interés del autor por el
tema, el contexto académico en el que se desarrolló y la utilidad que se espera
que tenga para el lector. Suele ocupar una o dos páginas.

Este trabajo nació en el marco del curso \emph{\monoCurso}, dictado en
\monoUniversidad, y responde a la necesidad de profundizar en un tema que, por
su relevancia, merece un tratamiento sistemático. Se ha procurado que el texto
sea claro, riguroso y verificable, de modo que pueda servir como punto de
partida para investigaciones posteriores.

Quiero expresar mi agradecimiento al docente \monoDocente, por su orientación;
a mis compañeros, por sus comentarios; y a mi familia, por su apoyo constante.
Cualquier error u omisión que subsista es de mi exclusiva responsabilidad.

\vspace{1cm}
\begin{flushright}
  \monoAutor\\
  \monoCiudad, \monoAnio
\end{flushright}
